\documentclass{article}
\usepackage[utf8]{inputenc}
\usepackage{xspace}
\usepackage[top=1.5cm, bottom=1.5cm, left=1.5cm, right=1.5cm]{geometry}
\usepackage{enumitem}
\usepackage{graphicx}
\usepackage{textcomp}
\usepackage{amssymb}
\usepackage{ifsym}
\usepackage{xspace}
\usepackage{amsmath}

\title{Facultad de Ciencias. Geoquímica orgánica}
\author{Arteaga Hernández Alejandro Iabin }
\date{Entrega 9 de octubre 2020}

\begin{document}

\maketitle

La geoquímica orgánica tuvo su origen en la primera parte del siglo XX, el padre de la geoquímica orgánica
es Alfred E Treibs.
\\

La geoquímica orgánica es la conjunción de dos ciencias distintas, la geología y la química orgánica,
a pesar de que el desarrollo de estas ciencias se dió de manera individual.
\\



El primer acercamiento entre estas dos ciencias fue cuando describió materia orgánica en sedimentos y 
rocas sedimentarias, otro acercamiento importante fue cuando se descubrió pigmentos orgánicos en el petróleo crudo y carbón, se demostró que estos pigmentos eran una degradación de substancias biológicas.
\\


\textbf{Parker D Trask}
Fue un geologo de la universidad de Berkeley, su campo de estudio eran rocas relacionadas al petróleo.
\\


\textbf{Alfred E Treibs}
Era un químico Alemán estudiado en Munich.
\\


La geoquímica orgánica es una rama interesante porque aglomera distintas ramas del conocimiento, 
desde la biología para el conocimiento de los procesos biológicos que pueden producir distintos materiales,
haciendo énfasis en el enfoque de la distribución de los hidrocarburos, en ciencias de la tierra desde el enfoque de la geología y estudio terreste, y la química que proporciona los métodos y técnicas de estudio
para el análisis de las distintas muestras.
\end{document}
